\section{Discussion}
\label{sec:discussion}

\subsection{Scale and clutter (RQ1)}
\begin{itemize}
    \item What made the tracker scale, and what it cost in accuracy.
    \item How effective the splitting is and why it fails: chaining of harbour vessels by new tracks' gates, truncation creating spurious dependence.
\end{itemize}

\subsection{Geometry (RQ2)}
\begin{itemize}
    \item No gain at this scale, as expected from Section~\ref{sec:nvector-error}. Why do it anyway: it is an investment for the thesis setting, at a cost of 1.3--1.7$\times$ run time.
    \item The rotation as a correction to the paper, independent of the results.
\end{itemize}

\subsection{Limits (RQ3)}
\begin{itemize}
    \item The long-revisit problem: gates grow, groups percolate into one, Murty dominates.
\end{itemize}

\subsection{Limitations of the study}
\begin{itemize}
    \item Seeds and confidence intervals; one region; mostly stationary vessels.
    \item Matched priors: birth rate and initial intensity taken from the truth.
    \item Constant $\PD$: tracks of vessels just outside the field of view count as false.
    \item Untuned process noise (position NEES 3.0, 2 if consistent).
    \item Simulated measurements on real tracks; no real sensor data.
\end{itemize}
